\paragraph{Statement A.4.4 (Sharp-arc exponent; complete candidate proof under review).}
For every $\epsilon>0$, the actual count in the question satisfies
\begin{equation}\label{A:sharp-two-thirds}
 R(X)=C_0X+O_\epsilon(X^{2/3+\epsilon}).
\end{equation}
The proof uses an angularly uniform Mellin transform. In particular, it
does not deduce an arc estimate from a count over an entire geodesic.

\begin{proof}
We give the transform and all uniform estimates used to invert it.
Put $A_z=(9z^2-1)^{1/2}$ and $\kappa_\ell=2\pi\ell/L_0$.
In this proof periods are unnormalized; this changes the estimates of
Statement A.4.3 only by fixed geodesic-length factors. Choose real
orthonormal eigenfunctions $u_j$ on $\Gamma^+\backslash\mathbb H$ and put
\[
 p_{j,\ell}=\int_{l_0}u_j(u)e^{i\kappa_\ell u}\,du,
 \qquad p_j'=\sum_{i=1}^2\int_{l_i}u_j\,du.
\]
Apart from the constant eigenfunction, write $s_j=1/2+it_j$, with
$t_j>0$, using Statement A.2.5. The constants in the following estimates
may depend on this fixed orbifold and its three geodesics. The spectral
estimates on its torsion-free cover apply to these eigenfunctions:
their lifts are normalized by the square root of the cover degree,
and a period on a lifted geodesic is the original period multiplied by
its fixed covering degree. The angular frequencies form the corresponding
fixed-index subsequence. Thus no spectral gap is assumed on the cover.

\emph{1. An exact Mellin identity.}
For $\operatorname{Re}s>1$ form the absolutely convergent series
\[
 E_s(p)=\sum_{i=1}^2\sum_{H_i\backslash\Gamma^+}
       \cosh\bigl(d(p,\gamma^{-1}l_i)\bigr)^{-s}.
\]
Absolute local uniform convergence follows by placing fixed disjoint
balls around coset representatives in the hyperbolic cylinder:
the number at distance at most $r$ is $O(e^r)$, since a cylinder
of fixed neck length has area $O(e^r)$ in that range. Its derivatives
also converge locally after increasing $\operatorname{Re}s$; the
identities initially obtained there continue to $\operatorname{Re}s>1$.
Set
\[
 K_s(\kappa)=\int_{\mathbb R}(\cosh u)^{-s}e^{i\kappa u}\,du
 =\frac{2^{s-1}\Gamma((s+i\kappa)/2)\Gamma((s-i\kappa)/2)}{\Gamma(s)},
\]
\[
 M_\ell(s)=K_s(\kappa_\ell)^{-1}
                 \int_{l_0}E_s(u)e^{i\kappa_\ell u}\,du.
\]
The beta integral proves the displayed formula for $K_s$.
For a disjoint target axis, the signed distance of a point on $l_0$
from that axis has absolute hyperbolic sine
$A\cosh(u-u_0)$, where $A=\sinh d(l_0,l)$ and $u_0$ is the
perpendicular foot. Consequently its contribution to $M_\ell$ is
\[
 e^{i\kappa_\ell u_0}A^{-s}F_{s,\kappa_\ell}(A),\quad
 F_{s,\kappa}(A)=
 {}_2F_1\left(\frac{s+i\kappa}{2},\frac{s-i\kappa}{2};
                    \frac{s+1}{2};-A^{-2}\right).
\]
To verify this identity without an unquoted transform formula, expand
$(1+A^{-2}\operatorname{sech}^2u)^{-s/2}$ and integrate termwise.
The coefficient is
\[
 \frac{(-1)^k(s/2)_k}{k!A^{2k}}
      \frac{K_{s+2k}(\kappa)}{K_s(\kappa)}
 =\frac{(-1)^k((s+i\kappa)/2)_k((s-i\kappa)/2)_k}
          {k!((s+1)/2)_k A^{2k}},
\]
by the gamma duplication formula. Here $A=A_z\geq\sqrt8$, so the
series and the original integrals converge absolutely.

There are four copies of each positive-height ideal angle: changing
the normal's sign and changing the sign of its $z$ coordinate leave
the foot coordinate unchanged modulo $L_0$. These are the four
regular cells already counted in Statement A.2.3. The two exceptional
double cosets with $z=0$ meet $l_0$ orthogonally at $u=L_0/2$.
Indeed $r_2=(1,1,0)$ is orthogonal to $e_z$ and its intersection
point has $\tanh u=1/\sqrt3$. Their individual contributions are
$e^{i\kappa_\ell L_0/2}=(-1)^\ell$. Thus exactly
\begin{equation}\label{A:Mellin-geometry}
 M_\ell(s)=2(-1)^\ell+4D_\ell(s)+4H_\ell(s),
\end{equation}
where
\[
 D_\ell(s)=\sum_{z\geq1}\sum_{N\mathfrak a=m_z}
          e^{2\pi i\ell\theta(\mathfrak a)}A_z^{-s},
\]
\[
 H_\ell(s)=\sum_{z\geq1}\sum_{N\mathfrak a=m_z}
 e^{2\pi i\ell\theta(\mathfrak a)}A_z^{-s}
                    (F_{s,\kappa_\ell}(A_z)-1).
\]
All sums in \eqref{A:Mellin-geometry} are absolutely convergent
in their initial half-plane. The half-turns reversing an axis do not
belong to its cyclic subgroup $H_i$; they therefore give the two
oriented normal cosets, as required for this factor of four.

Unfolding the scalar spectral coefficient of $E_s$ along $l_i$ gives
\[
 \left(\int_{l_i}u_j\right)
 \int_{\mathbb R}(1+x^2)^{-s/2}
 {}_2F_1\left(\frac{s_j}{2},\frac{1-s_j}{2};\frac12;-x^2\right)dx
 =\left(\int_{l_i}u_j\right)
 \frac{\sqrt\pi\Gamma((s+s_j-1)/2)\Gamma((s-s_j)/2)}{\Gamma(s/2)^2}.
\]
Here $x=\sinh r$ in Fermi coordinates. Integrating the eigenfunction
along the full target geodesic solves
$\xi''+\tanh r\,\xi'+\lambda_j\xi=0$; its odd initial-data solution
integrates to zero over $r\in\mathbb R$, and its even solution is
the displayed hypergeometric function. For the beta evaluation put
$y=x^2/(1+x^2)$, apply the power-series identity
\[
 {}_2F_1(s_j/2,(1-s_j)/2;1/2;-x^2)
 =(1+x^2)^{-s_j/2}{}_2F_1(s_j/2,s_j/2;1/2;y),
\]
and integrate the series against
$y^{-1/2}(1-y)^{(s+s_j-3)/2}$. The resulting Gauss beta sum
is the stated gamma quotient. All beta integrals are absolutely
convergent for $\operatorname{Re}s>1$.
It follows that, for $\ell\ne0$,
\begin{equation}\label{A:Mellin-spectral}
 M_\ell(s)=\sum_{j\geq1}p_j'p_{j,\ell}G(s,t_j,\kappa_\ell),
\end{equation}
\[
 G(s,t,\kappa)=
 \frac{\Gamma((s+1)/2)\Gamma((s+s_j-1)/2)\Gamma((s-s_j)/2)}
 {\Gamma(s/2)\Gamma((s+i\kappa)/2)\Gamma((s-i\kappa)/2)}.
\]
The constant eigenfunction vanishes because its nonzero Fourier
period on $l_0$ is zero. For each fixed $\ell$ the series in
\eqref{A:Mellin-spectral} converges normally on compact subsets away
from its poles, since its terms decay exponentially as
$t_j\to\infty$. It therefore supplies the needed meromorphic
continuation. In $\operatorname{Re}s>0$ its only possible poles are
$s=s_j$ and $s=1-s_j$.

\emph{2. A uniform bound for the geometric correction.}
Fix $0<\sigma<1/2$, write $s=\sigma+i\tau$, and put
$U=1+|\tau|$. With $\delta=\operatorname{arsinh}(1/A)$ one has
\begin{equation}\label{A:F-spherical-integral}
 F_{s,\kappa}(A)=
 \frac{\Gamma((s+1)/2)}{\sqrt\pi\Gamma(s/2)}
 \int_{-1}^1
 \bigl(\cosh(2\delta)+x\sinh(2\delta)\bigr)^{-(s+i\kappa)/2}
 (1-x^2)^{s/2-1}\,dx.
\end{equation}
For completeness, this identity follows by differentiating under
the integral and integrating once by parts in $x$: both sides,
as analytic even functions of $\delta$ near zero, solve
\[
 f''+2s\coth(2\delta)f'+(s^2+\kappa^2)f=0,
 \qquad f(0)=1,\quad f'(0)=0.
\]
Uniqueness follows coefficient by coefficient, since
$\operatorname{Re}s>0$; continuation along positive real $\delta$
then proves \eqref{A:F-spherical-integral}. Equivalently, it is
the quadratic hypergeometric transformation with parameters
$a+b=s$, $c=(s+1)/2$.
The base in the integral is positive. Its complex power has modulus
equal to the power $-\sigma/2$, independently of both $\tau$ and
$\kappa$; the weight has modulus $(1-x^2)^{\sigma/2-1}$.
Stirling's formula for the prefactor gives
\[
 |F_{s,\kappa}(A)|\ll_\sigma U^{1/2}.
\]
At $\delta=0$ the first derivative of the integrand is an odd
multiple of $x$, whose integral is zero. Its second derivative is
bounded by a constant times $(1+|s|+|\kappa|)^2$ times the same
integrable absolute weight, uniformly for
$0\leq\delta\leq\operatorname{arsinh}(1/\sqrt8)$.
Taylor's formula with integral remainder therefore gives
\begin{equation}\label{A:F-uniform}
 |F_{s,\kappa}(A)-1|
 \ll_\sigma U^{1/2}
       \min\{1,(U+|\kappa|)^2/A^2\}.
\end{equation}
Statement A.3.1 implies
$\#\{(z,\mathfrak a):A_z\leq V\}\ll V$.
Partial summation in \eqref{A:F-uniform} now proves normal
absolute convergence of $H_\ell(s)$ for $\operatorname{Re}s>0$
and, on our chosen line,
\begin{equation}\label{A:H-uniform}
 |H_\ell(\sigma+i\tau)|
       \ll_\sigma U^{1/2}(U+|\kappa_\ell|)^{1-\sigma}.
\end{equation}
The same convergence argument continues it to any closed strip
in $\operatorname{Re}s>-1$ avoiding the possible hypergeometric
parameter poles, but only the positive strip is needed here.

\emph{3. The remaining Mellin contour, uniformly in the angular cutoff.}
Let $T,H\geq2$. Let $w=w_T$ be a smooth function, equal to one
near zero, supported in $[0,2]$, with its transition confined to
$[1-1/T,1+1/T]$, $0\leq w\leq1$, and
$|w^{(k)}|\ll_kT^k$. Its Mellin transform satisfies, on every
fixed positive vertical strip and for every $N$,
\begin{equation}\label{A:w-Mellin-bound}
 |\widehat w(\sigma+i\tau)|
 \ll_{\sigma,N}U^{-1}(1+U/T)^{-N}.
\end{equation}
This follows by one integration by parts, whose derivative has
bounded $L^1$ norm, and repeated further integrations by parts in
the logarithmic variable. Let $b_\ell$, $1\leq\ell\leq H$, be
arbitrary complex numbers with $|b_\ell|\leq1/\ell$.
We prove
\begin{equation}\label{A:contour-uniform}
 \int_{\mathbb R}|\widehat w(\sigma+i\tau)|
 \left|\sum_{\ell=1}^H b_\ell M_\ell(\sigma+i\tau)\right|d\tau
       \ll_\sigma (T+H)\log(2H).
\end{equation}

Here are details of the gamma estimate and its summation, to make
the frequency uniformity explicit. Write
$p_{j,\ell}=a_{j,\ell}D_{s_j}(\kappa_\ell)$ as in
\eqref{A:Reznikov-mean}. Put
$a=3/4-\sigma/2$ and $b=1/2-\sigma/2$, so $1/2<a<3/4$.
For $\tau,t,\kappa\geq0$ Stirling's formula gives
\begin{align}\label{A:gamma-full-bound}
 |G(\sigma+i\tau,t,\kappa)D_{1/2+it}(\kappa)|
 &\ll_\sigma U^{1/2}
  \{(1+\tau+t)(1+|\tau-t|)\}^{-a}\\
 &\quad\cdot\{(1+\tau+\kappa)(1+|\tau-\kappa|)\}^{b}
 \{(1+t+\kappa)(1+|t-\kappa|)\}^{-1/4}
 e^{-c\operatorname{dist}(t,[\min(\tau,\kappa),\max(\tau,\kappa)])}.
 \nonumber
\end{align}
One may take $c=\pi/2$. The exponential follows exactly from
\[
 t-\max(t,\kappa)-\max(t,\tau)+\max(\tau,\kappa)
 =-\operatorname{dist}(t,[\min(\tau,\kappa),\max(\tau,\kappa)]).
\]
The same bound holds with $|\tau|$, $|\kappa|$ for arbitrary signs.
The fixed real parts of the numerator gamma functions are
$\sigma/2-1/4$, which are negative but not nonpositive integers;
thus the estimate is bounded at $\tau=\pm t$.

In a dyadic spectral band $J\leq1+t_j<2J$, the elementary
estimates and Reznikov's bound used in Statement A.4.3 give
\begin{equation}\label{A:three-basic-sums}
 \sum_j|p_j'|^2\ll J,\quad
 \sum_j|p_j'|\ll J^{3/2},\quad
 \sum_j|p_j'p_{j,\ell}|\ll J,
 \quad\sum_{|\ell|\leq K}|a_{j,\ell}|^2\ll\max(K,J).
\end{equation}
The middle bound uses the elementary $O(J^2)$ eigenvalue count.
Changing between $\ell$ and $\kappa_\ell$ changes dyadic bounds
only by fixed constants.

First suppose $\kappa_\ell\leq t_j/2$. Then
$|D_{s_j}(\kappa_\ell)|\asymp J^{-1/2}$.
For $\tau\geq t_j/2$, \eqref{A:gamma-full-bound} implies
\[
 |G(\sigma+i\tau,t_j,\kappa_\ell)|
 \ll_\sigma
 \left(\frac{1+\tau}{1+|\tau-t_j|}\right)^a
 e^{-c(t_j-\tau)_+}.
\]
For $\tau<t_j/2$ there is instead exponential decay in $J$,
with an immaterial polynomial factor. Since $a<1$, integration
over a dyadic $\tau$ block of size $U\gg J$, with the factor
$U^{-1}(1+U/T)^{-N}$, is
$O((1+U/T)^{-N})$; the same assertion holds for $U\asymp J$
because the integrable singular factor is bounded by
$(J/(1+|\tau-t_j|))^a$.
Summing \eqref{A:three-basic-sums} over $\ell$ costs $\log(2H)$.
Summing over $U\gg J$ and then dyadic $J$ costs $O(T)$:
for $J\leq T$ the contribution is
$O(J\log(2T/J))$, whose dyadic sum is $O(T)$, and the remaining
bands decay faster than any power of $J/T$.

Second suppose $\kappa_\ell\geq2t_j$, in an angular block
$K\leq\kappa_\ell<2K$ with $K\gg J$.
Up to $\tau\leq\max(2t_j,\kappa_\ell/2)$,
\eqref{A:gamma-full-bound} is bounded by
\[
 K^{1/2-\sigma}U^{1/2}J^{-a}(1+|\tau-t_j|)^{-a}
 e^{-c(\tau-t_j)_+}.
\]
In the remaining range, the exponential factor is at most
$e^{-c'K}$. For $\tau\gg K$ the polynomial factor is bounded
independently of $\tau$ (the total $\tau$ power in
\eqref{A:gamma-full-bound} is zero). For $\tau\asymp K$ it is
at most a fixed power of $K$. After the coefficient bounds below,
the sum of all these tails is therefore $O(\log(2T))$:
the integral of $U^{-1}(1+U/T)^{-N}$ above a fixed positive
lower limit is $O(\log(2T))$, and the exponential in $K$
absorbs the polynomial factors. This accounts also for
$\tau\geq\kappa_\ell$, where the exponential is
$e^{-c(\kappa_\ell-t_j)}$, rather than
$e^{-c(\tau-t_j)}$.
Cauchy--Schwarz in $\ell$, using
$\sum|a_{j,\ell}|^2\ll K$ and
$\sum_{\ell\asymp K}|b_\ell|^2\ll K^{-1}$, introduces no
additional factor. The elementary beta-integral bound
\[
 \int_0^\infty U^{-1/2}(1+|\tau-t_j|)^{-a}
 e^{-c(\tau-t_j)_+}(1+U/T)^{-N}\,d\tau
 \ll_\sigma J^{-a}\min(J,T)^{1/2}
\]
therefore bounds the spectral band by
\[
 K^{1/2-\sigma}J^{3/2-2a}\min(J,T)^{1/2}
 =K^{1/2-\sigma}J^\sigma\min(J,T)^{1/2}.
\]
Its dyadic sum over $J\ll K$ is $O_\sigma(K)$ for $K\leq T$
and $O_\sigma(\sqrt{KT})$ for $K>T$. Summing $K\ll H$ is
$O_\sigma(T+H)$.

Finally suppose $t_j/2<\kappa_\ell<2t_j$, so $J\ll H$.
Cauchy--Schwarz in $\ell$ gives
\[
 \left|\sum b_\ell a_{j,\ell}
       G(\sigma+i\tau,t_j,\kappa_\ell)D_{s_j}(\kappa_\ell)\right|
 \ll J^{-1/2}
       \left(\sum_{\kappa_\ell\asymp J}|GD|^2\right)^{1/2}.
\]
Squaring \eqref{A:gamma-full-bound} and using
$\sum_{\kappa_\ell\asymp J}(1+|t_j-\kappa_\ell|)^{-1/2}
\ll J^{1/2}$ gives, in the three respective $\tau$ ranges,
\[
 \sum_{\kappa_\ell\asymp J}|GD|^2\ll_\sigma
 \begin{cases}
 U/J,&\tau\leq t_j/2,\\
 J^{2b+1/2}(1+|\tau-t_j|)^{-2a},&t_j/2<\tau<2t_j,\\
 1,&\tau\geq2t_j.
 \end{cases}
\]
After multiplication by $|\widehat w|$, integration in $\tau$,
and summation of $|p_j'|$ over the spectral band, these give
respectively
\[
 O_\sigma(\sqrt{J\min(J,T)}),\qquad
 O_\sigma(J(1+J/T)^{-N}),\qquad
 O_\sigma\bigl(J\log(2+T/J)\min(1,(T/J)^N)\bigr).
\]
For the middle range the exponent is exactly
$b-1/4-a=-1/2$ before the final $J^{3/2}$ sum.
Their dyadic sum over $J\ll H$ is $O_\sigma(T+H)$.
These three cases prove \eqref{A:contour-uniform}. The finitely
many $0<t_j<1$ are included in the band $J=1$; the strict spectral
gap ensures all constants there are finite.

\emph{4. Shifting the contour and bounding the residues.}
For $Y\geq2$ write
\[
 W_\ell(Y)=\sum_{z\geq1}\sum_{N\mathfrak a=m_z}
 e^{2\pi i\ell\theta(\mathfrak a)}w(A_z/Y).
\]
Mellin inversion of \eqref{A:Mellin-geometry} on a line
$\operatorname{Re}s>1$ gives $4W_\ell$ as the inverse Mellin
transform of $M_\ell-2(-1)^\ell-4H_\ell$.
For $H_\ell$ shift to $\operatorname{Re}s=\sigma$ without
crossing a pole. Equations \eqref{A:H-uniform} and
\eqref{A:w-Mellin-bound} show, uniformly in the coefficients,
\begin{equation}\label{A:correction-inversion}
 \left|\sum_{\ell=1}^Hb_\ell\frac1{2\pi i}
 \int_{(\sigma)}\widehat w(s)Y^sH_\ell(s)\,ds\right|
 \ll_\sigma Y^\sigma
     T^{1/2}(T+H)^{1-\sigma}\log(2H).
\end{equation}
Indeed the remaining integral is bounded by
$\int_0^\infty U^{-1/2}(U+H)^{1-\sigma}(1+U/T)^{-N}\,d\tau$.
The exceptional constant $2(-1)^\ell$ contributes
$2(-1)^\ell w(1/Y)$ and hence $O(\log(2H))$ in the angular sum.

Shift the $M_\ell$ contour to the same line. Its residue at $s=s_j$ is
\[
 p_j'p_{j,\ell}\,r_j\frac{D_{s_j}(0)}{D_{s_j}(\kappa_\ell)},\qquad
 r_j=\frac{2\Gamma((s_j+1)/2)\Gamma(s_j-1/2)}{\Gamma(s_j/2)^3},
 \qquad |r_j|\ll(1+t_j)^{1/2}.
\]
The residue at $1-s_j$ has the conjugate gamma factors and obeys
the same estimate. Statement A.4.3 and
\eqref{A:w-Mellin-bound} therefore bound all residues, after summing
the angular coefficients, by
\begin{equation}\label{A:residues-inversion}
 \ll Y^{1/2}\log(2H)
      \sum_{J\text{ dyadic}}J^{1/2}(1+J/T)^{-N}
 \ll Y^{1/2}T^{1/2}\log(2H).
\end{equation}
The lower-half-plane residues use the same proof with conjugated
gamma ratios. No positivity or common phase of the periods is used.
The remaining contour costs
$O_\sigma(Y^\sigma(T+H)\log(2H))$ by
\eqref{A:contour-uniform}.

For clarity, all contour shifts can first be made at finite heights.
The eigenvalue bound $N(V)\ll V^2$ permits heights tending to infinity
whose distance from every $\pm t_j$ is at least a fixed multiple of
$V^{-3}$: the excluded intervals in $[V,V+1]$ have total length
$O(V^{-1})$. On these horizontal segments the gamma formula and
the spectral coefficient bounds give polynomial growth for each
fixed $H$. Taking $N$ arbitrarily large in
\eqref{A:w-Mellin-bound} makes their integrals tend to zero.
Absolute convergence of the residue series follows from
\eqref{A:residues-inversion}. The estimates on the resulting vertical
line are uniform in $H,T$, so this limiting argument makes no
uncontrolled exchange of a growing angular cutoff with a contour limit.

Combining these estimates proves the uniform smoothed inequality
\begin{equation}\label{A:smoothed-angular-sum}
 \left|\sum_{\ell=1}^H b_\ell W_\ell(Y)\right|
 \ll_\sigma\log(2H)\left[
 Y^{1/2}T^{1/2}+Y^\sigma(T+H)
 +Y^\sigma T^{1/2}(T+H)^{1-\sigma}+1\right].
\end{equation}
Because the coefficients $b_\ell$ are arbitrary within
$|b_\ell|\leq1/\ell$, choose their phases to align $W_\ell(Y)$.
Then the left side is exactly
$\sum_{\ell\leq H}|W_\ell(Y)|/\ell$.
This averaged angular estimate is sufficient; it avoids asserting
the stronger, individual-mode estimate \eqref{A:needed-twist}.

\emph{5. Returning to both sharp boundaries.}
The Erd\H{o}s--Tur\'an inequality used in Statement A.4.2 applies
also to a finite list with nonnegative real weights. For rational
weights this follows by replication after clearing denominators;
arbitrary weights follow by approximation and continuity. Hence the
weighted arc discrepancy is at most
\[
 O\left(\frac{W_0(Y)}H+
       \sum_{\ell\leq H}\frac{|W_\ell(Y)|}{\ell}+\log(2Y)\right).
\]
The endpoint term is $O(\log(2Y))$ by the Pell calculation already
proved; the other arc endpoint is empty.
Since Statement A.3.1 gives
$\#\{(z,\mathfrak a):A_z\leq V\}=(c/3)V+O(V^{2/3})$,
partial summation gives
\[
 W_0(Y)=\frac c3Y\widehat w(1)+O(Y^{2/3}),
 \qquad c=\frac{9L_0^2}{2\pi^2},
\]
uniformly for the monotone smooth cutoffs under discussion, because
their total variation is one.

Choose two such cutoffs $w^-,w^+$ with
$0\leq w^-\leq\mathbf1_{[0,1]}\leq w^+\leq1$,
transition intervals respectively $[1-1/T,1]$ and $[1,1+1/T]$.
They bracket the sharp radial count with nonnegative weights and
satisfy $\widehat w^\pm(1)=1+O(1/T)$.
Put $T=H=\lfloor Y^{1/3}\rfloor$ in
\eqref{A:smoothed-angular-sum}. The three nonconstant error terms
are respectively
\[
 O(Y^{2/3}\log Y),\quad
 O_\sigma(Y^{1/3+\sigma}\log Y),\quad
 O_\sigma(Y^{1/2+2\sigma/3}\log Y).
\]
The angular and radial main-term bracketing errors are
$O(Y/H+Y/T)=O(Y^{2/3})$.
Take any fixed sufficiently small $\sigma>0$, and absorb logarithms
into $Y^\epsilon$. Finally put $Y=(9X^2-1)^{1/2}$.
The ideal-arc main term is $(\lambda_0c/3)Y=C_0X+O(X^{-1})$.
Statement A.1.3 changes this by only $O(\log(2X))$, proving
\eqref{A:sharp-two-thirds}.
\end{proof}
